\documentclass[11pt,a4paper]{article}

\usepackage[utf8]{inputenc}
\usepackage[T1]{fontenc}
\usepackage[brazil]{babel}
\usepackage{amsmath}
\usepackage{amssymb}
\usepackage{booktabs}
\usepackage{geometry}
\usepackage{graphicx}
\usepackage{hyperref}
\usepackage{subcaption}
\usepackage{xcolor}
\usepackage{siunitx}

\geometry{margin=2.4cm}
\hypersetup{
    colorlinks=true,
    linkcolor=blue!50!black,
    urlcolor=blue!50!black,
    citecolor=blue!50!black
}

\title{Compara\c{c}\~ao Spyro--SPECFEM e an\'alises de converg\^encia\\
para propaga\c{c}\~ao de ondas el\'asticas 2D e 3D\\[6pt]
\large Valida\c{c}\~ao cruzada e casos com free surface no topo}
\author{}
\date{\today}

\begin{document}

\maketitle

% ============================================================================
\section{Objetivo}
% ============================================================================

Este relat\'orio documenta tr\^es grupos de resultados:
\begin{enumerate}
    \item \textbf{compara\c{c}\~ao direta Spyro--SPECFEM}: casos 2D e 3D
    usados como valida\c{c}\~ao cruzada, em configura\c{c}\~oes totalmente
    refletivas;
    \item \textbf{converg\^encia global por autoconverg\^encia}: erro $L^2$
    do campo de deslocamento no tempo final, medido contra uma solu\c{c}\~ao
    de refer\^encia em malha mais refinada (2D e 3D);
    \item \textbf{converg\^encia em receptores}: erro $L^2$ relativo dos
    tra\c{c}os nos receptores contra refer\^encia refinada (2D e 3D).
\end{enumerate}

\noindent
Todos os estudos de converg\^encia utilizam \emph{free surface} no topo
($z=0$) e condi\c{c}\~oes de Dirichlet nas demais faces.

% ============================================================================
\section{Escopo dos estudos}
% ============================================================================

Os resultados de converg\^encia cobrem dois cen\'arios, cada um avaliado com
elementos espectrais (quadril\'ateros/hex\'aedros, Q) e elementos KMV
(tri\^angulos/tetraedros, T):
\begin{enumerate}
    \item \textbf{global 2D}: dom\'inio $[0,\,3]\times[-3,\,0]$\,km, fonte em
    $(-0.6,\,1.5)$, Ricker $f=5$\,Hz com atraso $\tau=0.2$\,s,
    $t_f=1.0$\,s, $\Delta t=5\times10^{-4}$\,s. Malha de refer\^encia
    $h_{\text{ref}}=0.015$. O erro global \'e medido pela norma $L^2$
    ponderada do campo de deslocamento no tempo final, avaliada nos pontos
    de quadratura GLL da malha grosseira;
    \item \textbf{global 3D}: dom\'inio $[0,\,1]^3$\,km, fonte em
    $(-0.3,\,0.5,\,0.5)$, Ricker $f=5$\,Hz com atraso $\tau=0.2$\,s,
    $t_f=0.5$\,s, $\Delta t=5\times10^{-4}$\,s. Malha de refer\^encia
    $h_{\text{ref}}=0.04$.
\end{enumerate}

\noindent
Em todos os casos, a face superior em $z=0$ permanece livre (\emph{free
surface}) e as demais faces recebem condi\c{c}\~oes de Dirichlet
hom\^ogeneas. A compara\c{c}\~ao Spyro--SPECFEM \'e mantida como
refer\^encia de valida\c{c}\~ao cruzada e usa um setup separado, totalmente
refletivo.

% ============================================================================
\section{Compara\c{c}\~ao Spyro \emph{vs.}\ SPECFEM}
\label{sec:comparison}
% ============================================================================

A compara\c{c}\~ao direta com o SPECFEM \'e mantida neste relat\'orio como
evid\^encia de valida\c{c}\~ao cruzada do solver el\'astico do Spyro. Esses
casos usam um setup 2D/3D totalmente refletivo, separado dos estudos de
converg\^encia discutidos nas se\c{c}\~oes seguintes.

\subsection{Caso 2D}

No caso 2D, os dois solvers usam malhas equivalentes de grau 4 com
$h=0.02$, tempo final $t_f=1.5$\,s e passo temporal $\Delta t=10^{-3}$\,s. Os
erros L2 relativos nos receptores s\~ao:
\[
\varepsilon_{L2}(u_x) = 3.13\%, \qquad \varepsilon_{L2}(u_z) = 3.16\%.
\]

\begin{figure}[h!]
    \centering
    \begin{subfigure}{0.49\linewidth}
        \centering
        \includegraphics[width=\linewidth]{../results/elastic_evaluation/comparison_2d/ux_bxx_comparison.png}
        \caption{Componente $u_x$.}
    \end{subfigure}
    \hfill
    \begin{subfigure}{0.49\linewidth}
        \centering
        \includegraphics[width=\linewidth]{../results/elastic_evaluation/comparison_2d/uz_bxz_comparison.png}
        \caption{Componente $u_z$.}
    \end{subfigure}
    \caption{Compara\c{c}\~ao 2D Spyro \emph{vs.}\ SPECFEM2D: sismogramas nos receptores.}
    \label{fig:comp-2d}
\end{figure}

\begin{figure}[h!]
    \centering
    \includegraphics[width=\linewidth]{../results/elastic_evaluation/trace_detail_2d.png}
    \caption{Detalhe da compara\c{c}\~ao 2D: tra\c{c}os sobrepostos (topo), diferen\c{c}a pontual Spyro$-$SPECFEM (meio) e zoom na regi\~ao de maior erro (base).}
    \label{fig:trace-detail-2d}
\end{figure}

\subsection{Caso 3D}

No caso 3D, os dois solvers usam discretiza\c{c}\~oes equivalentes de grau 4
com $h=0.10$, tempo final $t_f=1.5$\,s e passo temporal $\Delta t=10^{-3}$\,s.
Os erros L2 relativos permanecem em torno de 3.6\%:
\[
\varepsilon_{L2}(u_x) = 3.66\%, \qquad
\varepsilon_{L2}(u_y) = 3.64\%, \qquad
\varepsilon_{L2}(u_z) = 3.63\%.
\]

\begin{figure}[h!]
    \centering
    \begin{subfigure}{0.49\linewidth}
        \centering
        \includegraphics[width=\linewidth]{../results/elastic_evaluation/comparison_3d/ux_bxx_comparison.png}
        \caption{Componente $u_x$.}
    \end{subfigure}
    \hfill
    \begin{subfigure}{0.49\linewidth}
        \centering
        \includegraphics[width=\linewidth]{../results/elastic_evaluation/comparison_3d/uy_bxy_comparison.png}
        \caption{Componente $u_y$.}
    \end{subfigure}
    \\[6pt]
    \begin{subfigure}{0.49\linewidth}
        \centering
        \includegraphics[width=\linewidth]{../results/elastic_evaluation/comparison_3d/uz_bxz_comparison.png}
        \caption{Componente $u_z$.}
    \end{subfigure}
    \hfill
    \begin{subfigure}{0.49\linewidth}
        \centering
        \includegraphics[width=\linewidth]{../results/elastic_evaluation/comparison_3d/trace_comparison.png}
        \caption{Tra\c{c}o do receptor central.}
    \end{subfigure}
    \caption{Compara\c{c}\~ao 3D Spyro \emph{vs.}\ SPECFEM3D.}
    \label{fig:comp-3d}
\end{figure}

\begin{figure}[h!]
    \centering
    \includegraphics[width=\linewidth]{../results/elastic_evaluation/trace_detail_3d.png}
    \caption{Detalhe da compara\c{c}\~ao 3D: tra\c{c}os sobrepostos, diferen\c{c}a pontual e zoom na regi\~ao de maior erro para cada componente ($u_x$, $u_y$, $u_z$).}
    \label{fig:trace-detail-3d}
\end{figure}

\noindent
Em ambos os casos, a concord\^ancia entre Spyro e SPECFEM \'e excelente, com
diferen\c{c}as residuais da ordem de 2--5\%, compat\'iveis com diferen\c{c}as de
implementa\c{c}\~ao na inje\c{c}\~ao da fonte, integra\c{c}\~ao temporal e
montagem da massa diagonal.

% ============================================================================
\section{Converg\^encia global por autoconverg\^encia}
\label{sec:global-selfconv}
% ============================================================================

Nesta se\c{c}\~ao, mede-se a taxa de converg\^encia do campo de deslocamento
completo no tempo final. O m\'etodo \'e o de \emph{autoconverg\^encia}:
executa-se uma simula\c{c}\~ao de refer\^encia em malha fina $h_{\text{ref}}$
e compara-se o campo de deslocamento das malhas mais grosseiras via norma
$L^2$ ponderada nos pontos de quadratura GLL. Para cada malha grosseira de
tamanho $h$, o campo de refer\^encia \'e interpolado na grade GLL da malha
grosseira usando um campo de refer\^encia tensor-produto Lagrangiano.

A fonte \'e um pulso de Ricker com frequ\^encia $f=5$\,Hz e atraso
$\tau=0.2$\,s. A escolha de $\tau=0.2$ garante que o wavelet \'e
essencialmente nulo em $t=0$ ($R(0)\approx -0.001$), evitando a
introdu\c{c}\~ao de descontinuidades iniciais que degradam a converg\^encia.

\subsection{Caso 2D}
\label{sec:global-2d}

Dom\'inio $[0,\,3]\times[-3,\,0]$\,km, grau $p=4$, malha de refer\^encia
$h_{\text{ref}}=0.015$ ($40\,000$ c\'elulas). Oito malhas de varredura
$h\in\{0.12, 0.10, 0.075, 0.06, 0.05, 0.04, 0.03, 0.02\}$. A ordem global
ajustada (regi\~ao $h\in[0.04,\,0.12]$) \'e:
\[
    p_{\text{global}} = 4.82 \qquad (\text{esperado}: 5.0).
\]

\begin{table}[h!]
\centering
\caption{Autoconverg\^encia 2D: erro $L^2$ relativo do campo de deslocamento.}
\label{tab:global-2d}
\begin{tabular}{@{}cccc@{}}
\toprule
$h$ & C\'elulas & $\varepsilon_{L^2}$ rel. & Taxa local \\
\midrule
0.120 &    625 & $9.54\times10^{-2}$ & --- \\
0.100 &    900 & $4.82\times10^{-2}$ & 3.75 \\
0.075 &  1\,600 & $6.71\times10^{-3}$ & 6.85 \\
0.060 &  2\,500 & $1.70\times10^{-3}$ & 6.14 \\
0.050 &  3\,600 & $4.64\times10^{-4}$ & 7.13 \\
0.040 &  5\,625 & $1.10\times10^{-4}$ & 6.46 \\
0.030 & 10\,000 & $6.73\times10^{-5}$ & 1.70 \\
0.020 & 22\,500 & $3.54\times10^{-5}$ & 1.59 \\
\bottomrule
\end{tabular}
\end{table}

\noindent
Para $h\le 0.04$, a taxa local cai para ${\approx}\,1.6$, indicando
satura\c{c}\~ao pelo erro da pr\'opria refer\^encia ($h/h_{\text{ref}}
\lesssim 2.7$). Nas malhas mais grosseiras ($h\in[0.04,\,0.12]$), a
converg\^encia \'e limpa e consistente com a ordem te\'orica.

\begin{figure}[h!]
    \centering
    \includegraphics[width=\linewidth]{../results/global_selfconv_top_free_2d/global_selfconv_top_free_2d.png}
    \caption{Autoconverg\^encia 2D: erro global do campo de deslocamento (esquerda) e erro nos receptores (direita). A reta-guia tracejada corresponde a $O(h^5)$.}
    \label{fig:global-2d}
\end{figure}

\subsection{Caso 3D}
\label{sec:global-3d}

Dom\'inio $[0,\,1]^3$\,km, grau $p=4$, malha de refer\^encia
$h_{\text{ref}}=0.04$ ($15\,625$ c\'elulas, ${\approx}\,3.1$\,M DOFs).
Quatro malhas de varredura $h\in\{0.125, 0.10, 0.0625, 0.05\}$. A ordem
global ajustada \'e:
\[
    p_{\text{global}} = 5.85 \qquad (\text{esperado}: 5.0).
\]

\begin{table}[h!]
\centering
\caption{Autoconverg\^encia 3D: erro $L^2$ relativo do campo de deslocamento.}
\label{tab:global-3d}
\begin{tabular}{@{}cccc@{}}
\toprule
$h$ & C\'elulas & $\varepsilon_{L^2}$ rel. & Taxa local \\
\midrule
0.1250 &    512 & $9.61\times10^{-2}$ & --- \\
0.1000 &  1\,000 & $3.55\times10^{-2}$ & 4.46 \\
0.0625 &  4\,096 & $1.46\times10^{-3}$ & 6.80 \\
0.0500 &  8\,000 & $5.67\times10^{-4}$ & 4.22 \\
\bottomrule
\end{tabular}
\end{table}

\noindent
A superestimativa da ordem ($5.85 > 5.0$) decorre das raz\~oes
$h/h_{\text{ref}}$ moderadas (3.1 a 1.25), t\'ipica do regime de
autoconverg\^encia quando a refer\^encia n\~ao \'e exata. Ainda assim, as
curvas por componente ($u_x$, $u_y$, $u_z$) seguem consistentemente a
reta-guia $O(h^5)$, confirmando que o solver 3D converge com a ordem
esperada no campo de deslocamento completo.

\begin{figure}[h!]
    \centering
    \includegraphics[width=\linewidth]{../results/global_selfconv_top_free_3d/global_selfconv_top_free_3d.png}
    \caption{Autoconverg\^encia 3D: erro global do campo de deslocamento (esquerda) e erro nos receptores (direita). A reta-guia tracejada corresponde a $O(h^5)$.}
    \label{fig:global-3d}
\end{figure}

% ============================================================================
\section{Resumo dos resultados de converg\^encia}
\label{sec:summary}
% ============================================================================

% ============================================================================
\subsection{Compara\c{c}\~ao Spectral \emph{vs.}\ KMV}
\label{sec:spec-vs-kmv}
% ============================================================================

Nesta se\c{c}\~ao, compara-se a converg\^encia dos elementos espectrais
(quadril\'ateros/hex\'aedros com mass lumping GLL, $p=4$) com os elementos
triangulares/tetra\'edricos KMV (Kong--Mulder--Veldhuizen, mass lumping em
simplices). No caso 2D, ambos usam grau $p=4$; no caso 3D, o KMV usa $p=3$
(m\'aximo suportado para tetraedros). Os tamanhos de malha $h$ do KMV
foram escolhidos de forma a produzir n\'umeros de graus de liberdade
similares aos correspondentes espectrais.

Para o KMV, o erro global \'e calculado por interpola\c{c}\~ao
\emph{cross-mesh} via Firedrake: ambas as solu\c{c}\~oes (refer\^encia em
malha fina e corrente em malha grosseira) s\~ao projetadas em um espa\c{c}o
DG intermedi\'ario na malha grosseira, permitindo o c\'omputo da norma
$L^2$ da diferen\c{c}a.

\subsubsection{Caso 2D}

\begin{figure}[h!]
    \centering
    \begin{subfigure}{0.49\linewidth}
        \centering
        \includegraphics[width=\linewidth]{../results/spectral_vs_kmv_2d_global_comparison.png}
        \caption{Erro global do campo.}
    \end{subfigure}
    \hfill
    \begin{subfigure}{0.49\linewidth}
        \centering
        \includegraphics[width=\linewidth]{../results/spectral_vs_kmv_2d_receiver_comparison.png}
        \caption{Erro nos receptores.}
    \end{subfigure}
    \caption{Converg\^encia 2D: Spectral ($p=4$) \emph{vs.}\ KMV ($p=4$).
    As retas tracejadas indicam a taxa te\'orica $O(h^5)$. Ambos os
    m\'etodos convergem com a ordem esperada na regi\~ao pr\'e-satura\c{c}\~ao.}
    \label{fig:spec-vs-kmv-2d}
\end{figure}

\noindent
A Tabela~\ref{tab:spec-vs-kmv-2d} resume as ordens observadas. Em 2D com
mesmo grau, o KMV apresenta erros absolutos ligeiramente maiores para o
mesmo $h$, compensados pelo fato de que cada c\'elula triangular possui
menos DOFs que um quadril\'atero de mesmo grau.

\begin{table}[h!]
\centering
\caption{Ordens observadas --- compara\c{c}\~ao 2D.}
\label{tab:spec-vs-kmv-2d}
\begin{tabular}{@{}lccc@{}}
\toprule
M\'etodo & Grau & Ordem esperada & Ordem observada (global) \\
\midrule
Spectral (Q) & 4 & 5.0 & 4.82 \\
KMV (T)      & 4 & 5.0 & 5.12 \\
\bottomrule
\end{tabular}
\end{table}

\subsubsection{Caso 3D}

No caso 3D, o KMV utiliza grau $p=3$ (m\'aximo suportado em tetraedros)
enquanto o spectral usa $p=4$ em hex\'aedros. A compara\c{c}\~ao \'e,
portanto, entre ordens te\'oricas distintas: $O(h^5)$ (spectral) e $O(h^4)$
(KMV).

\begin{figure}[h!]
    \centering
    \begin{subfigure}{0.49\linewidth}
        \centering
        \includegraphics[width=\linewidth]{../results/spectral_vs_kmv_3d_global_comparison.png}
        \caption{Erro global do campo.}
    \end{subfigure}
    \hfill
    \begin{subfigure}{0.49\linewidth}
        \centering
        \includegraphics[width=\linewidth]{../results/spectral_vs_kmv_3d_receiver_comparison.png}
        \caption{Erro nos receptores.}
    \end{subfigure}
    \caption{Converg\^encia 3D: Spectral ($p=4$) \emph{vs.}\ KMV ($p=3$).
    As retas tracejadas indicam $O(h^5)$ para spectral e $O(h^4)$ para KMV.}
    \label{fig:spec-vs-kmv-3d}
\end{figure}

\begin{table}[h!]
\centering
\caption{Ordens observadas --- compara\c{c}\~ao 3D.}
\label{tab:spec-vs-kmv-3d}
\begin{tabular}{@{}lccc@{}}
\toprule
M\'etodo & Grau & Ordem esperada & Ordem observada (global) \\
\midrule
Spectral (Q) & 4 & 5.0 & 5.85 \\
KMV (T)      & 3 & 4.0 & 3.55 \\
\bottomrule
\end{tabular}
\end{table}

\noindent
O KMV 3D converge com ordem ${\approx}\,3.55$, pr\'oximo do esperado
$O(h^4)$. A diferen\c{c}a de precis\~ao absoluta em rela\c{c}\~ao ao
spectral reflete tanto o grau inferior ($p=3$ \emph{vs.}\ $p=4$) quanto a
maior restri\c{c}\~ao de estabilidade CFL dos tetraedros (que exigiram
$\Delta t = 10^{-4}$\,s contra $5\times10^{-4}$\,s do spectral).

\subsection{Tabela geral}

\begin{table}[h!]
\centering
\caption{Principais resultados das an\'alises de converg\^encia.}
\label{tab:summary}
\begin{tabular}{@{}llclc@{}}
\toprule
M\'etodo & Dim. & Grau & M\'etrica & Ordem observada \\
\midrule
\multicolumn{5}{@{}l}{\textit{Erro global no campo de deslocamento (autoconverg\^encia)}} \\
Spectral (Q)  & 2D & 4 & $L^2$ ponderada, GLL      & 4.82 \\
Spectral (Q)  & 3D & 4 & $L^2$ ponderada, GLL      & 5.85 \\
KMV (T)       & 2D & 4 & $L^2$ cross-mesh DG        & 5.12 \\
KMV (T)       & 3D & 3 & $L^2$ cross-mesh DG        & 3.55 \\
\bottomrule
\end{tabular}
\end{table}

\noindent
As principais conclus\~oes s\~ao:
\begin{enumerate}
    \item Na m\'etrica global (campo completo), a ordem te\'orica $p+1$ \'e
    atingida ou ligeiramente superada em todos os casos. Para o spectral,
    $p+1=5$ \'e observado tanto em 2D ($4.82$) quanto em 3D ($5.85$). Para o
    KMV, observa-se $5.12$ em 2D (esperado 5) e $3.55$ em 3D (esperado 4).
    \item A leve superestimativa no spectral 3D decorre do regime de
    autoconverg\^encia, com raz\~oes $h/h_{\text{ref}}$ moderadas.
    \item O KMV 3D utiliza grau $p=3$ (m\'aximo para tetraedros) e apresenta
    restri\c{c}\~oes de CFL mais severas ($\Delta t = 10^{-4}$\,s contra
    $5\times10^{-4}$\,s do spectral), resultando em maior custo
    computacional por passo temporal.
\end{enumerate}

% ============================================================================
\section*{Refer\^encia}
% ============================================================================

\noindent
[1] Brenner, S. C.; Scott, L. R. \textit{The Mathematical Theory of Finite
Element Methods}. 3rd ed. New York: Springer, 2008.

\end{document}
